\documentclass[a4paper, 12pt]{report}

\usepackage[utf8]{inputenc}
\usepackage{color}
\usepackage{geometry}
\usepackage[pdftex]{graphicx}
\usepackage{tikz}

\usepackage{filecontents}
\usepackage[T1]{fontenc}
\usepackage[UKenglish]{babel}
\usepackage{newpxtext,newpxmath}
\usepackage[babel=true]{csquotes}
\usepackage[round]{natbib}


% Logos
\newcommand{\ulb}{\includegraphics[scale=1.1]{PageDeGarde_MFE/logo_ULB2.pdf}}
\newcommand{\polytech}{\includegraphics[scale=0.35]{PageDeGarde_MFE/logo_polytech_FR.pdf}}

% Polices
\definecolor{ULBblue}{rgb}{0,0.2196,0.5765}
\newcommand{\fontTitle}{\sffamily \Huge\selectfont \color{ULBblue}}
\newcommand{\fontSubtitle}{\sffamily \LARGE \selectfont \color{ULBblue}}
\newcommand{\fontText}{\sffamily \selectfont}
\newcommand{\fontColor}{\sffamily \selectfont \color{ULBblue}}

% Titre
\newcommand{\titleA}{\fontTitle{Human - Robots Swarms Interactions}} % Titre identique au titre remis au secrétariat
%\newcommand{\titleB}{\fontTitle{Deuxième ligne de titre du mémoire}} % (dans la langue de rédaction a priori)
% Sous-titre
\newcommand{\subtitle}{\fontSubtitle{Ligne du sous-titre du mémoire}}
% Titre du diplôme
\newcommand{\diplomaA}{\fontText{Mémoire présenté en vue de l’obtention du diplôme}} % A laisser en Français
\newcommand{\diplomaB}{\fontText{d'Ingénieur Civil en Informatique à finalité Intelligence Computationnelle}}

% Etudiant
\newcommand{\student}{\textbf{\sffamily \large Anthony Debruyn}}

% Supervision
\newcommand{\promAa}{\fontColor{Directeur}}
\newcommand{\promAb}{\fontText{Professeur [Prénom Nom]}}
\newcommand{\promBa}{\fontColor{Co-Promoteur}}
\newcommand{\promBb}{\fontText{Professeur [Prénom Nom]}}
\newcommand{\promCa}{\fontColor{Superviseur}}
\newcommand{\promCb}{\fontText{[Prénom Nom]}}
\newcommand{\deptA}{\fontColor{Service}}
\newcommand{\deptB}{\fontText{[Nom du service]}}

% Année académique
\newcommand{\yearA}{\fontColor{Année académique}}
\newcommand{\yearB}{\fontText{2014 - 2015}}

\newcommand{\quoto}[2]{
\begin{quotation}
\begin{it}
\enquote{#1} - #2
\end{it}
\end{quotation}
}

\newcommand{\quot}[1]{\textit{\enquote{#1}}}

\begin{document}

	\thispagestyle{empty}
	\newgeometry{top=2.5cm, bottom=1.5cm, left=2.5cm, right=1cm}
	\setlength{\unitlength}{1mm}
	\noindent\begin{picture}(175,257)
	
		\put(0,245){\polytech}
		\put(153,139.5){\ulb}
		
		\put(8,155){\makebox(150,10)[l]{\titleA}}
%		\put(8,145){\makebox(150,10)[l]{\titleB}}
		\put(8,135){\makebox(150,10)[l]{\subtitle}}
		
		\put(0,75){
		\begin{tikzpicture}[scale=0.1]
		\fill [fill=ULBblue](0,0) rectangle (0.8,90);
		\fill [fill=ULBblue](0,57) rectangle (152,57.8);
		\end{tikzpicture}}
		
		\put(8,120){\makebox(150,5)[l]{\diplomaA}}
		\put(8,115){\makebox(150,5)[l]{\diplomaB}}
		
		\put(8,75){\makebox(150,10)[l]{\selectfont \student}}
		
		\put(8,44){\makebox(80,5)[l]{\promAa}}
		\put(8,39){\makebox(80,5)[l]{\promAb}}
		\put(8,31){\makebox(80,5)[l]{\promBa}} % Commenter la ligne si pas nécessaire
		\put(8,26){\makebox(80,5)[l]{\promBb}} % Commenter la ligne si pas nécessaire
		\put(8,18){\makebox(80,5)[l]{\promCa}} % Commenter la ligne si pas nécessaire
		\put(8,13){\makebox(80,5)[l]{\promCb}} % Commenter la ligne si pas nécessaire
		\put(8,5){\makebox(80,5)[l]{\deptA}}
		\put(8,0){\makebox(80,5)[l]{\deptB}}
		
		\put(145,5){\makebox(30,5)[r]{\yearA}}
		\put(145,0){\makebox(30,5)[r]{\yearB}}
	
	\end{picture}
	\restoregeometry
	
% Template conçu par Benjamin Vanhemelryck et revu par François Bronchart - Mai 2013
	
%%%%%%%%%%%%%%%%%%%%%%%%%%%%%%%%%%%%%%%%%%%%%%%%%%%%%%%
% TABLE OF CONTENT
%%%%%%%%%%%%%%%%%%%%%%%%%%%%%%%%%%%%%%%%%%%%%%%%%%%%%%%

\tableofcontents
	
%%%%%%%%%%%%%%%%%%%%%%%%%%%%%%%%%%%%%%%%%%%%%%%%%%%%%%%
% INTRODUCTION
%%%%%%%%%%%%%%%%%%%%%%%%%%%%%%%%%%%%%%%%%%%%%%%%%%%%%%%

\chapter{Introduction}

[To write at the end. When I have a global view.]

\chapter{The Problem}
	\section{Introduction}
	
	In this section, we will discuss the problem that led to the creation of this thesis by first providing the reader with some general insight in the world of swarm robotics and swarm intelligence. Then we will focus on specific parts of these domains of study: design methods, feedbacks between human and robots' swarm. Eventually, we will conclude with the main problem this thesis is trying to solve.
	
	\section{Swarm Engineering}
	
	This section is largely inspired by~\citet{brambilla2013swarm}, a reviewing article on swarm engineering. For \citet{csahin2005swarm}, swarm robotics is defined as \quot{the study of how large numbers of relatively simple physically embodied agents can be designed such that a desired collective behavior emerges from the local interactions among agents and between the agents and the environment}~\citep{csahin2005swarm}. Swarm robotics can be separated from other robotic studies by the following characteristics \citep{brambilla2013swarm}:
	
\begin{itemize}
\item Robots are \emph{autonomous}
\item Robots evolve \emph{in the environment} and can interact with it
\item Robots' interactions are \emph{local} (sensors and communications)
\item No \emph{centralised control} or \emph{global knowledge}
\item Robots \emph{cooperate} to achieve a certain goal
\end{itemize}

As in this field of study, one is always looking for \emph{robust}, \emph{scalable} and \emph{flexible} systems, the main source of inspiration is the group of social animals: ants, birds, fishes, ... When some of these simple animals gather in groups, they are able to perform tasks that could not be achieved individually (collective behaviour emerges from local interactions). Below are listed the definitions of these three terms \citep{brambilla2013swarm}:

\begin{description}
\item[Robustness:] Resistance against \emph{loss of group entities}. One can increase it by adding redundancy or remove the need for a leader.
\item[Scalability:] Low variation in the performance of a system with respect to the \emph{size of the system}. It can be increased by encouraging local interactions, such as sensing and communications.
\item[Flexibility:] Low variation in the performance of a system with respect to the \emph{type of environment or the task}.
\end{description}

With these definitions in mind, we can explain swarm engineering as:

\quoto{Swarm engineering is an emerging discipline that aims at defining systematic and well founded procedures for modeling, designing, realizing, verifying, validating, operating, and maintaining a swarm robotics system.}{\cite{brambilla2013swarm}}

\citet{kazadi2000swarm} points out that \quot{to the swarm engineer, the important points in the design of a swarm are that the swarm will do precisely what it is designed to do, and that it will do so reliably and on time} \citep{kazadi2000swarm}.
	
	\section{Design Methods}
	
	\citet{brambilla2013swarm} did a pretty good job in classifying the different methods to design and 
	\section{Feedbacks}
	\section{Robots - Human Feedback}
	
[Currently, only little research on feedback between human and robots swarms. That research is focused on interaction with an additional layer... Mostly unidirectional communication (human to robots). Need new types of interactions for new applications. Robots to human (guide). Write about article saying self-organised feedback is better.]



	\section{Protection and Exploration}
	
	The final objective of this thesis is to create a protective behaviour for a swarm of robots
	

\chapter{Solution}

	\section{Introduction}
	\section{Swarm Robotics}
	\section{The Shape}
	\section{Virtual Physics}
	\section{Potentials}
		\subsection{Human Potential}
		\subsection{Gravity Potential}
		\subsection{Agent Repulsion Potential}
		\subsection{Default Potential}
	\section{Range and Bearing Sensor}
		\subsection{Tests}
		\subsection{Calibration}
		\subsection{Conclusions}

	\section{Omnidirectional Camera Sensor}
		\subsection{Tests}

	\section{The Hardware}
		\subsection{Objective}
		\subsection{Choices}
		\subsection{Blueprints}
		\subsection{Build Process}

\chapter{Assessments}
	\section{Metrics}
	\section{More robots ?}
	\section{Tests Scenarios}
	\section{Human Travelling Speed ?}
	\section{Usability Study}

\chapter{Future Works}
	\section{Other Robots}
	\section{Human Guidance}
		\subsection{Zero Visibility Areas or Blind People}
		\subsection{Human Motion Synchronisation}
		\subsection{Vehicle Guidance}

\chapter{Acknowledgements}

[Never gonna give you up\\
Never gonna let you down\\
Never gonna run around and desert you\\
Never gonna make you cry\\
Never gonna say goodbye\\
Never gonna tell a lie and hurt you]

\newpage

%%%%%%%%%%%%%%%%%%%%%%%%%%%%%%%%%%%%%%%%%%%%%%%%%%%%%%%
% BIBLIOGRAPHY
%%%%%%%%%%%%%%%%%%%%%%%%%%%%%%%%%%%%%%%%%%%%%%%%%%%%%%%

\begin{filecontents}{thesis.bib}
%http://en.wikibooks.org/wiki/LaTeX/Bibliography_Management

@inproceedings{podevijn2012self,
  title={Self-organised feedback in human swarm interaction},
  author={Podevijn, Ga{\"e}tan and O’Grady, Rehan and Dorigo, Marco},
  booktitle={Proceedings of the workshop on robot feedback in human-robot interaction: how to make a robot readable for a human interaction partner (Ro-Man 2012)},
  year={2012}
}

@incollection{podevijn2014gesturing,
  title={Gesturing at subswarms: Towards direct human control of robot swarms},
  author={Podevijn, Ga{\"e}tan and O’Grady, Rehan and Nashed, Youssef SG and Dorigo, Marco},
  booktitle={Towards Autonomous Robotic Systems},
  pages={390--403},
  year={2014},
  publisher={Springer}
}

@article{brambilla2013swarm,
  title={Swarm robotics: a review from the swarm engineering perspective},
  author={Brambilla, Manuele and Ferrante, Eliseo and Birattari, Mauro and Dorigo, Marco},
  journal={Swarm Intelligence},
  volume={7},
  number={1},
  pages={1--41},
  year={2013},
  publisher={Springer}
}

@incollection{csahin2005swarm,
  title={Swarm robotics: From sources of inspiration to domains of application},
  author={{\c{S}}ahin, Erol},
  booktitle={Swarm robotics},
  pages={10--20},
  year={2005},
  publisher={Springer}
}

@phdthesis{kazadi2000swarm,
  title={Swarm engineering},
  author={Kazadi, Sanza T},
  year={2000},
  school={California Institute of Technology}
}

\end{filecontents}

\nocite{*}
\bibliographystyle{plainnat}
\bibliography{thesis}

\end{document}
